\section{Discussion}
\label{sec:discussion}

\begin{outlinenote}
\begin{enumerate}[leftmargin=*,nosep]
\item Opening, no heading, two sentences: the one major finding . Users expect deletion to mean gone and cannot check it, and what happened depended on the agent, not on the wording.
\item Subsection: What Explains the Differences Across Agents. Where each agent keeps the fact and which data deletion method reaches that place; what each agent's documentation says a method removes, set against what we observed; regulation (we looked for a difference and report what we found, or that none distinguishes the agents). Explanations are worded as consistent with the observations, because the audit is black-box.
\item Subsection: Why Users Cannot Tell. The check users reach for first misleads; the agent's reply and the interface do not show the outcome. This joins the survey (expectation, check) to the audit. [Heading wording open.]
\item Subsection: What Each Party Should Do, by audience. Platforms: state what each data deletion method removes, say what a deletion did and did not do, and give users a way to check. Regulators: deletion rights have to be testable (no claim of GDPR compliance or non-compliance). Users: which methods reached what we tested, marked provisional. Researchers: the audit method and data are available. A table follows (as in Khandelwal et al.), with three columns: what we found, what we recommend, who acts. Each row traces a recommendation to a result section.
\item Subsection: Limitations, short and last, no new heading level. The survey describes US Prolific users of the six agents; free-tier accounts and one time window; prompt-based recall is weak evidence that a fact is gone; the can't tell share (provisional until recall finishes); the audit reports what a user can observe, not what a platform stores. The free-text reasons participants gave (not coded) go here as a note without counts.
\item Left out: repeated numbers (Sections 4 and 5 hold them), future work (Conclusion), related work (its own section).
\end{enumerate}
\end{outlinenote}
